\documentclass[12pt,a4paper]{article}

\usepackage[T1]{fontenc}
\usepackage[utf8]{inputenc}
\usepackage[brazil]{babel}
\usepackage{lmodern}
\usepackage{geometry}
\usepackage{hyperref}
\usepackage{enumitem}
\usepackage{array}

\geometry{margin=2.5cm}
\hypersetup{
  colorlinks=true,
  linkcolor=blue,
  urlcolor=blue
}

\title{Plano de Aula: Mapeamento do Modelo ER para o Modelo Relacional}
\date{}

\begin{document}
\maketitle

\noindent
\textbf{Duração:} 2 horas (120 minutos)\\
\textbf{Público-alvo:} Iniciantes\\
\textbf{Recursos Necessários:} Quadro branco/negro, marcadores/giz, folhas/cartões, post-its, (opcional) projetor.

\section{Objetivos da Aula}
\begin{itemize}[leftmargin=*,itemsep=0.25em]
  \item Gerar esquemas relacionais a partir de um esquema conceitual (ER).
  \item Aplicar regras de transformação para entidades, atributos e relacionamentos.
  \item Diferenciar mapeamento de relacionamentos 1:1, 1:N e N:M.
  \item Identificar PKs e FKs resultantes do mapeamento.
\end{itemize}

\section{Metodologia}
Aula demonstrativa com regras de mapeamento no quadro e prática em grupos transformando um ER (produzido na Aula 2) em tabelas. A correção é feita comparando soluções e justificando escolhas (onde colocar FKs, como tratar N:M).

\section{Cronograma e Conteúdo Detalhado (120 Minutos)}
\subsection*{Parte 1: Revisão do ER (15 minutos)}
\begin{itemize}[leftmargin=*,itemsep=0.35em]
  \item Revisar rapidamente: entidades, atributos, relacionamentos e cardinalidade.
  \item Definir o ER base que será mapeado (ex.: Biblioteca).
\end{itemize}

\subsection*{Parte 2: Regras de Mapeamento (35 minutos)}
\begin{itemize}[leftmargin=*,itemsep=0.35em]
  \item \textbf{Entidades viram tabelas (10 min)}
  \begin{itemize}[leftmargin=*,itemsep=0.25em]
    \item Atributos simples viram colunas.
    \item Identificador vira PK.
  \end{itemize}
  \item \textbf{Relacionamento 1:N (10 min)}
  \begin{itemize}[leftmargin=*,itemsep=0.25em]
    \item FK do lado 1 vai para o lado N.
    \item Analogia: ``todo PET aponta para um CLIENTE''.
  \end{itemize}
  \item \textbf{Relacionamento 1:1 (5 min)}
  \begin{itemize}[leftmargin=*,itemsep=0.25em]
    \item Estratégias: FK em uma tabela ou mesclar (dependendo do caso).
  \end{itemize}
  \item \textbf{Relacionamento N:M (10 min)}
  \begin{itemize}[leftmargin=*,itemsep=0.25em]
    \item Criar tabela associativa com duas FKs.
    \item PK composta (ou ID próprio).
  \end{itemize}
\end{itemize}

\subsection*{Parte 3: Exemplo Guiado (20 minutos)}
\begin{itemize}[leftmargin=*,itemsep=0.35em]
  \item Mapear no quadro um relacionamento N:M (Aluno--Disciplina).
  \item Mostrar tabela associativa (MATRICULA) com FKs e atributos do relacionamento (ex.: data\_matricula).
\end{itemize}

\subsection*{Parte 4: Atividade em Grupo --- ER \textrightarrow\ Tabelas (40 minutos)}
\begin{itemize}[leftmargin=*,itemsep=0.35em]
  \item \textbf{Entrada:} um diagrama ER (Biblioteca ou Clínica) com pelo menos 3 entidades.
  \item \textbf{Tarefas:}
  \begin{itemize}[leftmargin=*,itemsep=0.25em]
    \item Transformar entidades em tabelas com PK.
    \item Aplicar regras para cada relacionamento (1:1, 1:N, N:M).
    \item Identificar FKs e justificar onde elas ficaram.
  \end{itemize}
  \item \textbf{Entrega:} esquema relacional (lista de tabelas com PK/FK).
\end{itemize}

\subsection*{Parte 5: Encerramento (10 minutos)}
\begin{itemize}[leftmargin=*,itemsep=0.35em]
  \item Revisar regras principais e erros comuns (FK no lado errado, esquecer tabela associativa).
  \item Ponte: ``Próxima aula: operações do modelo relacional e como consultar os dados.''
\end{itemize}

\end{document}

